\begin{sang}{Der kan man C}{Morten Schou, 2026 - Mel: Røvervise (Vi lister os afsted på tå)}
\spal2
\begin{vers}
Jeg finder datamaten frem
når jeg skal programmere.
Det er lidt svært, men bli'r først slemt,
når det skal kompilere.
Den kaster op og skaber sig.
Jeg tænker ik' det rager mig
om min parameter er const eller ej,
eller navne må skrives med binde-streg.
\end{vers}
\begin{vers}
Jeg lærer lidt om modulo,
forstår dog ikke resten.
Bli'r jeg en rigtig datalog?
Består jeg unit testen?
Min nye ven ChatGPT
fortæller mig, enhver kan C,
at pointeren peger i ring på sig selv -
induktionsalgoritmen bli'r triviel.
\end{vers}
\begin{vers}
Det går nok, hvis jeg skriver if,
men hvem er hend' der Else?
Kan nogen give mig et fif?
Betingelser er trælse.
De vigtigste får !udråbstegn
men så' der noget andet i vej'n.
Jeg bryder mig ik' om programmer med breaks,
nu kan switchene udføre hver en case.
\end{vers}
\begin{vers}
Programmet itererer ej
min while er foruløkket,
Kan rekursionen redde mig?
Er linjen galt indrykket?
Funktionen spørger kun sig selv,
men fortsætter alligevel.
Den hopper af kæden og lister afsted.
Nu' der rod i mit træ uden fejlbesked.
\end{vers}
\begin{vers}
Min aflevering blev for lang,
har glemt nul-terminering.
Jeg havde ellers denne gang
brugt tid på kommentering.
Jeg caster typerne i flæng,
men ASCII teksten er for streng.
De siger det blot er en pointer til char,
allokeret med malloc på heap - hurra!??
\end{vers}
\vbox{}\vfill
\vbox{}\vfill
\vbox{}\vfill
\vbox{}\vfill
\vbox{}\vfill
\vbox{}\vfill
\vbox{}\vfill
\laps
\end{sang}